% Standalone problem document, generated from the adjacent README.md.
% Compile with XeLaTeX. Mathematical content is maintained in Markdown.
\documentclass[11pt,a4paper]{article}
\usepackage[fontsize=10.5pt]{scrextend}
\usepackage[margin=22mm,headheight=15pt,headsep=9mm,footskip=11mm]{geometry}
\usepackage{fontspec}
\setmainfont{texgyrepagella}[Extension=.otf,UprightFont=*-regular,BoldFont=*-bold,ItalicFont=*-italic,BoldItalicFont=*-bolditalic]
\setsansfont{texgyreheros}[Extension=.otf,UprightFont=*-regular,BoldFont=*-bold,ItalicFont=*-italic,BoldItalicFont=*-bolditalic]
\setmonofont{texgyrecursor}[Extension=.otf,UprightFont=*-regular,BoldFont=*-bold,ItalicFont=*-italic,BoldItalicFont=*-bolditalic,Scale=MatchLowercase]
\usepackage{amsmath,amssymb}
\usepackage{unicode-math}
\setmathfont{texgyrepagella-math.otf}
\usepackage{xcolor,microtype,enumitem,needspace,fancyhdr}
\usepackage{bookmark}
\definecolor{ink}{HTML}{17344B}
\definecolor{accent}{HTML}{197D80}
\definecolor{muted}{HTML}{596773}
\hypersetup{colorlinks=true,linkcolor=accent,urlcolor=accent,pdftitle={TR-04 - Improve
the worst-case approximation factor for prescribed tensor-train
ranks},pdfauthor={Open Problems in Numerical Linear Algebra}}
\urlstyle{same}
\setlength{\parindent}{0pt}
\setlength{\parskip}{5pt plus 1pt minus 1pt}
\linespread{1.04}
\setlength{\emergencystretch}{3em}
\widowpenalty=10000
\clubpenalty=10000
\displaywidowpenalty=10000
\allowdisplaybreaks[1]
\setlist{nosep,leftmargin=1.5em}
\providecommand{\tightlist}{\setlength{\itemsep}{0pt}\setlength{\parskip}{0pt}}
\providecommand{\pandocbounded}[1]{#1}
\pagestyle{fancy}
\fancyhf{}
\fancyhead[L]{\sffamily\footnotesize\color{muted}OPEN PROBLEMS IN NUMERICAL LINEAR ALGEBRA}
\fancyhead[R]{\sffamily\bfseries\footnotesize\color{accent}TR-04}
\fancyfoot[L]{\parbox[b]{0.9\textwidth}{\sffamily\fontsize{8}{10}\selectfont\color{muted}Editorial ratings; a bounded search cannot certify that a problem remains unsolved.\\Literature check: 2026-09-08}}
\fancyfoot[R]{\sffamily\footnotesize\color{muted}\thepage}
\renewcommand{\headrulewidth}{0.3pt}
\renewcommand{\headrule}{\hbox to\headwidth{\color{accent}\leaders\hrule height \headrulewidth\hfill}}
\makeatletter
\renewcommand{\section}{\Needspace{5\baselineskip}\@startsection{section}{1}{0pt}{12pt plus 2pt minus 2pt}{4pt}{\sffamily\bfseries\large\color{ink}}}
\renewcommand{\subsection}{\Needspace{4\baselineskip}\@startsection{subsection}{2}{0pt}{9pt plus 1pt minus 1pt}{3pt}{\sffamily\bfseries\normalsize\color{ink}}}
\makeatother
\setcounter{secnumdepth}{0}
\begin{document}
{\sffamily\small\color{muted}Tensor computations\par}
\vspace{3pt}
{\raggedright\hyphenpenalty=10000\exhyphenpenalty=10000\sffamily\bfseries\fontsize{21}{24}\selectfont\color{ink}Improve
the worst-case approximation factor for prescribed tensor-train
ranks\par}
\vspace{7pt}
{\sffamily\small\textbf{Difficulty:} extreme\quad\textbf{Importance:} interesting
to the community\par}
\vspace{4pt}
{\color{accent}\hrule height 0.8pt}
\vspace{6pt}
\textbf{Status:} open; explicit source updated 2026-08-20.

\subsection{Problem statement}\label{problem-statement}

Let \(d\ge3\), \(n_1,\ldots,n_d\ge2\), and positive integers
\(r_1,\ldots,r_{d-1}\) be given. Let
\(S_r\subset\mathbb R^{n_1\times\cdots\times n_d}\) consist of tensors
whose unfolding separating modes \(1,\ldots,j\) from modes
\(j+1,\ldots,d\) has rank at most \(r_j\), for every \(j\). For a dense
input tensor \(A\), put

\[
E_*(A,r)=\min_{Y\in S_r}\|A-Y\|_F^2.
\]

Find a polynomial-time algorithm returning \(X\in S_r\) with

\[
\|A-X\|_F^2<(d-1)E_*(A,r)
\qquad\text{whenever }E_*(A,r)>0,
\]

and exact reconstruction when \(E_*(A,r)=0\); alternatively, establish a
complexity obstruction to such a guarantee. Ranks may not be increased.
Polynomial time is measured in the dense input size and rank parameters,
in the idealized arithmetic/SVD model used for TT-SVD; a bit-complexity
formulation must additionally specify precision and output tolerances.

This is the tensor-train case of the published tree-network question.
The positive-error qualification corrects the impossible strict
inequality at zero optimum. It asks for a uniformly valid algorithm, not
an empirical improvement or a guarantee restricted to particular input
tensors.

\subsection{References}\label{references}

I. V. Oseledets,
\href{https://doi.org/10.1137/090752286}{\emph{Tensor-Train
Decomposition}}, \emph{SIAM Journal on Scientific Computing} 33 (2011),
2295--2317, Theorem 2.2 and Corollary 2.4. Amsel et al.,
\href{https://arxiv.org/html/2602.05394v3}{workshop report}, §6.1,
Problem 6.1. Matthew Fahrbach and Mehrdad Ghadiri,
\href{https://arxiv.org/html/2508.06693v1}{\emph{A Tight Lower Bound for
the Approximation Guarantee of Higher-Order Singular Value
Decomposition}}, Theorems 1.2--1.3, concerns tightness of specific
Tucker algorithms rather than a lower bound against all tensor-train
algorithms.

\subsection{Status check}\label{status-check}

Searches included
\textit{"tensor train" "approximation" "2026" "improvement"} and
\textit{"tree tensor" "approximation guarantee" "2026"}. No
universal improvement or matching complexity lower bound was located.
Randomized methods and better practical error estimates require a
separate comparison with the displayed worst-case statement.
\end{document}
